% !PRINTED 462
\weberpage{462}
Ist $h_1 = 2$, $h_2 = h_3 = h_4 = 0$, so ist $m = 12$. Eine
Invariantencurve $12^{\text{ter}}$ Ordnung $\Phi$ muss die Curve $f$ in
den 24 Punkten $P_7$ berühren, und wenn wir in $\Phi - a\Delta^2$ die
Constante $a$ passend bestimmen, so ergiebt sich, dass diese Function
durch $f$ theilbar sein muss. Der Quotient kann, als Invariante
$8^{\text{ter}}$ Ordnung, nur von der Form $bf^2$ sein, und demnach ist
$\Phi$ von der Form $a\Delta^2 + bf^3$. Ebenso können wir schliessen,
dass eine Invariante $18^{\text{ter}}$ Ordnung die Form
$a\Delta^3 + b\Delta f^3$ haben muss.

\begin{center}
\gesperrt{Alle Invarianten $6^{\text{ter}}$, $12^{\text{ter}}$, $18^{\text{ter}}$
Ordnung sind also rationale Functionen von $\Delta$ und $f$.}
\end{center}

Der nächste Werth, den $m$ nach (1) haben kann, ist $m = 14$. In
diesem Falle ist $h_2 = 1$, während $h_1, h_3, h_4$ gleich Null sind.
Es gehen also alle Invarianten $14^{\text{ter}}$ Ordnung durch die 56
Punkte $P_3$, und diese bilden das vollständige Schnittpunktsystem
einer solchen Invariantencurve mit der Grundcurve. Dies gilt auch von
der Invariante $C$. Haben wir eine zweite Invariante $14^{\text{ter}}$
Ordnung $C'$, so können wir wieder, wie oben, die Constante $a$ so
bestimmen, dass $C' - aC$ durch $f$ theilbar ist. Der Quotient ist
eine Invariante $10^{\text{ter}}$ Ordnung, und daraus folgt, da es
ausser $f\Delta$ keine Invariante $10^{\text{ter}}$ Ordnung giebt:

\begin{center}
\gesperrt{Jede Invariante $14^{\text{ter}}$ Ordnung ist in der Form
darstellbar:}
\end{center}
\[
aC + bf^2\Delta,
\]
\begin{center}
\gesperrt{worin $a, b$ Constanten sind. Umgekehrt ist jeder Ausdruck von
dieser Form eine Invariante $14^{\text{ter}}$ Ordnung.}
\end{center}

Es kann sodann $m$ nach (1) den Werth 20 haben, nämlich für
$h_1 = h_2 = 1$. Eine Invariante $20^{\text{ster}}$ Ordnung muss also
durch die Punkte $P_7$ und $P_3$ gehen, d.\,h. durch die 80
Schnittpunkte von $f$ mit $\Delta C$. Daraus können wir ebenso wie
vorhin schliessen, dass eine Invariante $20^{\text{sten}}$ Grades in
der Form
\[
f(a\Delta^2 + bf^3) + cC\Delta.
\]
darstellbar ist, worin $a, b, c$ beliebige Constanten sind. Eine
unabhängige Invariante $20^{\text{ster}}$ Ordnung giebt es nicht.

Nehmen wir nun an, es seien $K'$ und $K$ zwei Invarianten
$21^{\text{ster}}$ Ordnung; beide müssen nach (1) durch die 84 Punkte
$\overline{P}_2$ gehen, und folglich kann man $a$ so bestimmen, dass
$K' - aK$ durch $f$ theilbar wird. Der Quotient wäre eine Invariante
%CONT>>
